\documentclass[10pt,a4paper]{amsart}
\usepackage[margin=19mm]{geometry}
\usepackage{amsmath,amssymb,mathtools}
\usepackage[scr=rsfs,cal=euler]{mathalfa}
\usepackage{xfrac}
\usepackage[unicode,hidelinks,bookmarks=false]{hyperref}
\newcommand{\ie}{i.e.\ }
\newcommand{\cf}{cf.\ }
\newcommand{\resp}{respectively\ }
\newcommand{\Subsection}{\subsection}
\providecommand{\nobreakdash}{\nobreak\hbox{-}}
\setlength{\emergencystretch}{2em}
\multlinegap0pt
\usepackage{amssymb}
\usepackage{enumerate}
\usepackage{xcolor}
\newtheorem{theorem}{Theorem}
\newtheorem{claim}{Claim}
\title{\textbf{A somewhat sure note on an un-Schur problem}}
\author{Swaroop Hegde, Hitesh Kumar, Pratibha}
\date{Source: arXiv:2609.18474v1 (CC BY 4.0)}
\begin{document}
\maketitle

\noindent\emph{Source and licence.} \textbf{A somewhat sure note on an un-Schur problem}. Version arXiv:2609.18474v1 (CC BY 4.0), distributed by arXiv under the Creative Commons Attribution 4.0 International licence, http://creativecommons.org/licenses/by/4.0/, as recorded in the arXiv licence metadata for this exact version and shown on the abstract page https://arxiv.org/abs/2609.18474v1. Full licence notice: \texttt{licenses/arxiv-2609.18474-CC-BY-4.0.txt}.\par\smallskip

\noindent\textit{Editorial scope.} Counted unit: the article's theorem thm:lower\_bound\_3\_color (source0.tex lines 189--200). The article's complete original proof of it is delivered with it (source0.tex lines 201--306, one \texttt{proof} environment of 106 lines). No mathematics is written, repaired or re-derived: the statement and the proof are the article's own lines, in the article's own notation, and no retained line is substituted or re-flowed.  No further source-local context is needed: every cross-reference of every retained line resolves inside the two delivered spans.  Nothing was omitted from any retained span, so the package carries no omission record.  No retained line contains a citation, so the wrapper carries no bibliography and no bibliography entry.  No retained line contains a figure environment, an image-inclusion command or drawing code, so no image is delivered and the package declares no asset dependency.  Editorial formatting only, in addition: an AMS \texttt{amsart} preamble that restates the article's own declarations and macros used by the retained lines, namely the environments \texttt{theorem}, \texttt{claim} on separate counters, together with the standard packages those lines need.  The retained environments print in this document as Theorem 1, Claim 1 in order of appearance, whereas the article prints the same items with the numbering of its own full document.  The complete native source of the article as distributed in the arXiv e-print for this exact version is bundled unchanged as \texttt{sources/210701/source0.tex}.
\begin{theorem}\label{thm:lower_bound_3_color}
There exists a sequence of colorings $c_n:[n]\to [3]$ for which
\[
|R_n(c_n)|
=
\left(\frac{9}{44}+o(1)\right)n^2.
\]
Consequently,
\[
\liminf_{n\to \infty} \Lambda_{n,3}\ge \frac{9}{22}.
\]
\end{theorem}

\begin{proof}
For $n\in \mathbb{N}$, take
\begin{align*}
A& =\left\{i\in[n]: i\text{ is odd and }i\leq\frac{4n}{11}\right\},\\
B& =\left\{i\in[n]: i\text{ is odd and }
\frac{4n}{11}<i\leq\frac{10n}{11}\right\},\\
C&=\left\{i\in[n]: i\text{ is odd and }i>\frac{10n}{11}\right\}.
\end{align*}
Clearly, $A\sqcup B\sqcup C$ is a partition of all odd integers in $[n]$. Define the coloring $c_n$ as follows:
\[
c_n(i)=
\begin{cases}
1 & i\in A;\\
2 & i\in B;\\
1 & i\in C;\\
3 &  i \text{ is even}.
\end{cases}
\]

Next, we count the rainbow Schur triples $(x,y,x+y)$ w.r.t. the coloring $c_n$ according to the parity of the summands $x,y$. By the definition of $c_n$, both $x$ and $y$ cannot be even, or else they receive the same color. So we can partition
\[ R_n(c_n) = R_{\mathrm{oo}} \sqcup R_{\mathrm{oe}},\]
where $R_{oo}$ contains the pairs $(x,y)$ where both $x,y$ are odd, and $R_{\mathrm{oe}}$ contains the pairs where exactly one of $x,y$ is odd. 

\begin{claim}\label{claim:R_oe} We have 
\[|R_{\mathrm{oe}}| = \left(\frac{15}{121}+o(1)\right)n^2.\]
\end{claim}

\begin{proof}
Let $(x,y)\in R_{\mathrm{oe}}$. Then $x+y$ is odd. Assume first that $x$ is odd. Then $y$ is even and $c_n(y)=3$. Also, $c_n(x), c_n(x+y)\in \{1,2\}$ and $c_n(x)\neq c_n(x+y)$. It is then clear that 
\[ (x,x+y)\in A\times B \quad \text{or} \quad (x,x+y)\in B\times C.\]
Thus, the number of such triples is $|A||B| + |B||C|$. By symmetry between $x$ and $y$, it follows that 
\begin{align*}
    |R_{\mathrm{oe}}| & = 2\left(|A||B| + |B||C|\right)\\
    & = 2\left(\frac{2}{11}\frac{3}{11}+\frac{3}{11}\frac{1}{22}\right)n^2+o(n^2)\\
    &=\left(\frac{15}{121}+o(1)\right)n^2.\qedhere
\end{align*}
\end{proof}

\begin{claim}\label{claim:R_oo} We have 
\[|R_{\mathrm{oo}}| = \left(\frac{39}{484}+o(1)\right)n^2.\]
\end{claim}

\begin{proof} Let $(x,y)\in R_{\mathrm{oo}}$. Then $x+y$ is even, and hence $c_n(x+y) = 3$. Also, $c_n(x), c_n(y)\in \{1,2\}$ and $c_n(x)\neq c_n(y)$. Clearly, $x,y$ both cannot belong to $A\cup C$. Moreover, if $x\in B$ and $y\in C$, then
\[
x+y>
\frac{4n}{11}+\frac{10n}{11}>n.
\]
We conclude that $(x,y)$ cannot have one element in $B$ and the other in $C$. It remains to consider pairs $(x,y)$ having one element in $A$ and one in $B$. Define
\[
N_{AB}
=
\left\{(x,y)\in A\times B:x+y\leq n\right\} \quad \text{and}\quad N_{BA}
=
\left\{(x,y)\in B\times A:x+y\leq n\right\}.
\]
It is clear by symmetry that $|N_{AB}| = |N_{BA}|$, and
\[R_{\mathrm{oo}} = N_{AB}\sqcup N_{BA}.\]
So to estimate $|R_{oo}|$, it suffices to estimate $|N_{AB}|$. Fix an odd $y\in B$. The number of odd $x\in A$ satisfying
$x+y\leq n$ is
\[
\frac12
\min\left\{\frac{4n}{11},\,n-y\right\}+O(1).
\]
Using a Riemann-sum approximation, we obtain
\[
|N_{AB}|
=
\frac{n^2}{4}
\int_{4/11}^{10/11}
\min\left\{\frac{4}{11},1-t\right\}\,dt
+O(n).
\]
The point at which the two terms inside the minimum coincide is
$t=7/11$. Hence,
\[
\begin{aligned}
\int_{4/11}^{10/11}
\min\left\{\frac{4}{11},1-t\right\}\,dt
&=
\int_{4/11}^{7/11}\frac{4}{11}\,dt
+
\int_{7/11}^{10/11}(1-t)\,dt\\
&=
\frac{12}{121}+\frac{15}{242}\\
&=
\frac{39}{242}.
\end{aligned}
\]
It follows that
\[
|N_{AB}|
=
\left(\frac{39}{968}+o(1)\right)n^2,
\]
and therefore
\[
|R_{\mathrm{oo}}|
=
2|N_{AB}|
=
\left(\frac{39}{484}+o(1)\right)n^2.\qedhere
\]
\end{proof}
Using Claims \ref{claim:R_oe} and \ref{claim:R_oo}, we conclude that
\[    |R_n(c_n)| = |R_{\mathrm{oe}}|+|R_{\mathrm{oo}}| = \left(\frac{9}{44}+o(1)\right)n^2.\qedhere\]
\end{proof}

\end{document}
